% =========================================================================
\section{Circuito 2: amplificador no inversor}\label{sec:pl-no-inversor}
% =========================================================================
\begin{prelab}[Amplificador no inversor]\label{pl:no-inversor}
Usando un modelo de amplificador operacional ideal, encuéntrese la ganancia en tensión
$A_V=\Vout/\Vin$ en función de las resistencias del circuito de la \cref{fig:no-inversor}. Si
imponemos $R_1=1\,\mathrm{k}\Omega$, determínese el valor de $R_2$ para conseguir $A_V=+11$.
\end{prelab}

% Figura: amplificador no inversor (circuito 2) con tensiones y corrientes del análisis.
\begin{figure}[ht]
\centering
\begin{circuitikz}[circuito]
  \draw (0,0) node[op amp, ao, noinv input up] (A) {};
  % nudo inversor (abajo) y red de realimentación
  \draw (A.out) -- ++(0.7,0) coordinate (O);
  \draw (A.-) -- ++(-0.6,0) coordinate (N) -- ++(0,-1.0) coordinate (M);
  \draw (M) -- ++(-1.0,0) coordinate (G1) to[R, l_=$R_1$, i>^=$i_1$] ++(0,-1.7) node[ground] {};
  \draw (M) to[R, l_=$R_2$, i<_=$i_2$] (M -| O) -- (O);
  % generador en la entrada no inversora
  \draw (A.+) to[short, -o] ++(-2.0,0) coordinate (Ein) node[above, tens] {$\vin$};
  \draw (Ein) -- ++(-0.8,0) to[sV, l_=$v_g$] ++(0,-2.0) node[ground] {};
  % carga y salida
  \draw (O) to[short, *-o] ++(1.0,0) node[right, tens] {$\vout$};
  \draw (O) ++(0.5,0) to[R, l=$R_L$] ++(0,-2.0) node[ground] {};
  \fill (M) circle (1.6pt);
  \node[tens, anchor=east, font=\scriptsize] at ($(N)!0.5!(M)+(-0.08,0)$) {$\vm=\vin$};
\end{circuitikz}
\caption{Amplificador no inversor (figura~3 del guion~\cite{GUION-S1}), con $R_1=1\,\mathrm{k}\Omega$,
$R_L=10\,\mathrm{k}\Omega$ y la alimentación omitida. La entrada $+$ no absorbe corriente; el nudo
inversor sigue a la entrada ($\vm=\vin$) y $R_1$--$R_2$ forman un divisor de $\vout$.}
\label{fig:no-inversor}
\end{figure}


\begin{resolucion}
Con los sentidos de corriente de la \cref{fig:no-inversor}:

\textbf{1. Cortocircuito virtual.} La señal entra por $+$ y, por (R1), no hay caída en esa conexión;
por (R2),
\begin{equation}\label{eq:ni-vm}
  \vm=\vp=\vin .
\end{equation}

\textbf{2. Ley de Ohm en $R_1$} (del nudo inversor a tierra):
\begin{equation}\label{eq:ni-i1}
  i_1=\frac{\vm-0}{R_1}=\frac{\vin}{R_1} .
\end{equation}

\textbf{3. Kirchhoff en el nudo inversor}, con $i_-=0$ por (R1):
\begin{equation}\label{eq:ni-kcl}
  i_2=i_1+i_-=\frac{\vin}{R_1} .
\end{equation}

\textbf{4. Ley de Ohm en $R_2$} ($i_2$ va de la salida al nudo inversor):
\begin{equation}\label{eq:ni-vout}
  \vout=\vm+i_2R_2=\Bigl(1+\frac{R_2}{R_1}\Bigr)\vin .
\end{equation}

\textbf{Resultado.}
\begin{equation}\label{eq:ni-Av}
  \resultado{A_V=\frac{\vout}{\vin}=1+\frac{R_2}{R_1}}
\end{equation}

\textbf{Diseño.} Con $A_V=11$ y $R_1=1\,\mathrm{k}\Omega$ en \cref{eq:ni-Av}:
\begin{equation}\label{eq:ni-R2}
  R_2=(A_V-1)\,R_1\quad\Longrightarrow\quad\resultado{R_2=10\,\mathrm{k}\Omega}
\end{equation}

\textbf{Comentario.} $A_V>0$: salida en fase con la entrada, y $A_V\ge1$ siempre. Como en el inversor,
la ganancia solo depende del cociente de resistencias (\cref{fig:ni-graficas}).
\end{resolucion}

% Figura: previsiones del amplificador no inversor (A_V=+11, Vsat=9 V, f=1 kHz).
\begin{figure}[ht]
\centering
\begin{tikzpicture}
  \begin{axis}[tema, name=a, width=0.45\textwidth, height=4.6cm,
      xmin=0, xmax=2, ymin=-10.5, ymax=10.5, ytick={-9,-5.5,0,5.5,9},
      xlabel={$t$ (ms)}, ylabel={Tensión (V)}, title style={font=\small},
      title={(a) $V_p=0{,}5\,\mathrm{V}$}]
    \addplot[asintota, domain=0:2, samples=2] {9};
    \addplot[asintota, domain=0:2, samples=2] {-9};
    \addplot[senal entrada, domain=0:2, samples=150] {0.5*sin(360*x)};
    \addplot[senal salida, domain=0:2, samples=150] {5.5*sin(360*x)};
  \end{axis}
  \begin{axis}[temacentro, at={(a.east)}, anchor=west, xshift=1.2cm,
      width=0.42\textwidth, height=4.6cm,
      xmin=-1.6, xmax=1.6, ymin=-11, ymax=11, xtick={-0.82,0.82}, ytick={-9,9},
      xticklabels={$-0{,}82$,$0{,}82$}, yticklabels={$-9$,$9$},
      xlabel={$\vin$ (V)}, 
      xlabel style={at={(axis description cs:1,0.5)}, anchor=south east, yshift=2pt, font=\scriptsize}, clip=false,
      title style={font=\small, yshift=10pt}, title={(b) $\vout(\vin)$}]
    \addplot[senal salida, domain=-1.6:1.6, samples=200] {max(-9,min(9,11*x))};
    \addplot[punteado] coordinates {(0.818,0) (0.818,9)};
    \addplot[punteado] coordinates {(-0.818,0) (-0.818,-9)};
    \node[font=\scriptsize, anchor=west, text=rojoteo] at (axis cs:-1.55,5.8) {pendiente $+11$};
    \node[font=\scriptsize, anchor=west] at (axis description cs:0.52,0.97) {$\vout$ (V)};
  \end{axis}
\end{tikzpicture}
\caption{Previsiones para el amplificador no inversor ($A_V=+11$, $\Vsat\approx9\,\mathrm{V}$,
$f=1\,\mathrm{kHz}$). (a) Entrada (azul) y salida (ámbar) en fase. (b) Característica de
transferencia: pendiente $+11$ hasta $\abs{\vin}=\Vsat/11\approx0{,}82\,\mathrm{V}$ y saturación fuera.}
\label{fig:ni-graficas}
\end{figure}

